\begin{afterword}
\summaryhead

The post starts from a reported pattern: groups hit by a failed prophecy or a scandal about their founder often come back with stronger belief. The usual explanation is cognitive dissonance in each member, who cannot admit a costly mistake.

The author proposes a second mechanism, modeled on evaporative cooling, in which the fastest atoms escape a trap and the rest grow colder. After a shock, the most skeptical members leave first. The debate among those who remain then runs from fanatics to slightly less extreme fanatics, and without its moderates the group may grow more fanatical still. The illustrations are the member who walked out in \textsc{When Prophecy Fails}, the Objectivists who left after the Rand--Branden scandal, and a mailing list whose libertarians were insulted into leaving.

The lesson is to be slow to eject dissenters. The post ends with a counterweight: Kuhn's view that a science must shut out outsiders before it can do real work, and a footnote on excluding trolls from online discussions.

\responsehead

The idea is plausible, and the post's main source bears it out in part. \textsc{When Prophecy Fails} reports that the member who left on the night of the failed prophecy was the group's most vocal skeptic, and that the only two members who gave up the belief were lightly committed.

The post contrasts the idea with a narrower reading of that source than the book supports. It presents the cognitive-dissonance account as purely individual. Festinger and his colleagues made social support one of their five conditions, and themselves reported that the least committed left. What the post adds is the suggestion that these departures change the members who stay.

That suggestion is the part left unsupported. That a group's average moves when its moderates leave is arithmetic, and needs no change of mind in anyone. That those who remain ``may internally increase in average fanaticism'' is hedged, and none of the three examples tests it. Each shows who left.

The examples are also loosely reported. The Jehovah's Witnesses and Unarius show groups surviving a failed date, and the Witnesses lost active members after 1975. The Objectivist example joins two splits about twenty years apart. The affair was not made public in 1968, and the ``open system'' faction formed around David Kelley in 1989, not around Branden.

The effect the post sets out to explain is itself disputed. Melton argued that testing strengthens religious groups. Kelly argues that failed prophecies are often fatal, and, from newly unsealed archives, that the book's own group dissolved and its leader recanted. Even the book reports that the group won no convert and was gone within three weeks.

Finally, the last paragraph pulls against the advice before it. Kuhn, cited as the ``flip side,'' describes sciences that read dissenters ``out of the profession'' and then work more effectively. The post does not say how to tell that exclusion from the kind it warns against.

In short: a plausible mechanism, partly anticipated by the study it cites, illustrated by examples that show who leaves a group but not what happens to those who stay, and followed by a counterweight from Kuhn that the post does not reconcile with its advice.
\end{afterword}
